% Representative shape of the beam-model mixture density for one ray.
%
%   Sum of the four parts of the beam model, drawn with illustrative values
%   (not the rover's parameters):
%     p_short : exponential decay on [0, z*), cut off at the expected range z*
%     p_rand  : uniform floor over [0, r_max]
%     p_hit   : Gaussian centered on the expected range z*
%     p_max   : narrow box at r_max
%
% Pure vector TikZ so it survives the latex -> dvipdf build.
\begin{tikzpicture}[>=Latex, line join=round, font=\footnotesize, x=1.6cm, y=2cm,
    density/.style={black, line width=0.8pt},
    axis/.style={->, black!85, line width=0.6pt},
    guide/.style={black!40, dashed, line width=0.4pt},
  ]
  \def\zs{3.2}   % expected range z*
  \def\rmx{5.5}  % maximum range r_max

  % axes
  \draw[axis] (0,0) -- (6.0,0) node[below] {$z$};
  \draw[axis] (0,0) -- (0,2.0) node[left] {$p(z \mid x_t, m)$};

  % p_rand + p_short + p_hit on [0, r_max)
  \draw[density] plot[domain=0:\rmx-0.08, samples=300, variable=\x]
    ({\x}, {0.12 + (\x<\zs)*0.45*exp(-0.5*\x) + 1.5*exp(-((\x-\zs)^2)/0.0648)});

  % p_max: narrow box at r_max on top of the uniform floor
  \draw[density] (\rmx-0.08,0.12) -- (\rmx-0.08,1.0) -- (\rmx,1.0) -- (\rmx,0);

  % tick labels
  \draw[guide] (\zs,0) -- (\zs,1.62);
  \draw (\zs,0) -- (\zs,-0.05) node[below] {$z_t^{i*}$};
  \draw (\rmx,0) -- (\rmx,-0.05) node[below] {$r_{\max}$};
\end{tikzpicture}

%%% Local Variables:
%%% mode: latex
%%% TeX-master: "../thesis.tex"
%%% End:
